\section{Thirteen things happen first}
\label{sec:ch03-pipeline}
\index{request pipeline|(}

Mosaic logs this when it starts:

\begin{minted}[fontsize=\footnotesize]{text}
Request pipeline: 13 middleware
  1. InstrumentationMiddleware
  2. ExceptionMiddleware
  3. TimeoutMiddleware
  4. DocumentCacheMiddleware
  5. DocumentParserMiddleware
  6. DocumentValidationMiddleware
  7. CostAnalyzerMiddleware
  8. OperationCacheMiddleware
  9. OperationResolverMiddleware
  10. SkipWarmupExecutionMiddleware
  11. OperationVariableCoercionMiddleware
  12. ConcurrencyGateMiddleware
  13. OperationExecutionMiddleware
\end{minted}

\index{request pipeline!default}%
Twelve of those are HotChocolate's default pipeline, and they are added by a
method that is refreshingly easy to read:

\begin{minted}{csharp}
internal static void AddDefaultPipeline(this IList<RequestMiddlewareConfiguration> pipeline)
{
    pipeline.Add(CommonMiddleware.Instrumentation);
    pipeline.Add(CommonMiddleware.UnhandledExceptions);
    pipeline.Add(TimeoutMiddleware.Create());
    pipeline.Add(CommonMiddleware.DocumentCache);
    pipeline.Add(CommonMiddleware.DocumentParser);
    pipeline.Add(CommonMiddleware.DocumentValidation);
    pipeline.Add(OperationCacheMiddleware.Create());
    pipeline.Add(OperationResolverMiddleware.Create());
    pipeline.Add(CommonMiddleware.SkipWarmupExecution);
    pipeline.Add(OperationVariableCoercionMiddleware.Create());
    pipeline.Add(CommonMiddleware.ConcurrencyGate);
    pipeline.Add(OperationExecutionMiddleware.Create());
}
\end{minted}

That is the whole of it, in \code{RequestExecutorBuilderExtensions.UseRequest.cs}
at tag 16.6.0. There is no configuration system behind it and no discovery: the
default pipeline is a list of twelve entries appended in order.

\index{cost analysis!middleware}%
The thirteenth is the cost analyzer, which is not in that list because it
inserts itself. The defaults turn it on unless you pass
\code{disableDefaultSecurity}, and \code{AddCostAnalyzer} ends by asking to be
put in a specific place:

\begin{minted}{csharp}
.UseRequest(
    CostAnalyzerMiddleware.Create(),
    after: WellKnownRequestMiddleware.DocumentValidationMiddleware)
\end{minted}

Which is why it lands at seven, between validation and the operation cache.
Chapter~\ref{ch:hotchocolate-quickly} met the cost analyzer as the thing that
sprinkles \code{@cost} directives into exported SDL; this is the rest of it,
doing work on every request.

\subsection*{Named neighbours, and why they exist}

\index{request pipeline!middleware keys}%
That \code{after:} argument is worth dwelling on, because it explains the
strings in the log. Every middleware in the pipeline carries an optional key,
and the framework's own keys are public constants on
\code{WellKnownRequestMiddleware}. So an insertion does not have to count
positions:

\begin{minted}[fontsize=\footnotesize]{csharp}
builder.AddGraphQL()
    .UseRequest(
        MyMiddleware.Create(),
        after: WellKnownRequestMiddleware.DocumentValidationMiddleware);
\end{minted}

The key of the thing being inserted travels inside the configuration that
\code{Create()} returns, which is why it does not appear in that call. There are
overloads that take a bare delegate and a separate \code{key} argument, and one
that takes a middleware type; they all end up in the same list.

If the anchor is missing the call throws, naming the key it could not find,
which is a better failure than silently landing in the wrong place. Position
counting would break the first time a package inserted something above you, and
packages do: this is exactly how the cost analyzer gets in.

\subsection*{The pipeline is built once, not per request}

\index{request pipeline!composition}%
The list becomes a delegate chain in \code{CreatePipeline}, and the shape of
that method is the thing to take away:

\begin{minted}{csharp}
if (pipeline.Count == 0)
{
    defaultPipelineFactory ??= RequestExecutorBuilderExtensions.AddDefaultPipeline;
    defaultPipelineFactory(pipeline);
}

foreach (var modifier in pipelineModifiers)
{
    modifier(pipeline);
}

RequestDelegate next = _ => default;

for (var i = pipeline.Count - 1; i >= 0; i--)
{
    next = pipeline[i].Middleware(factoryContext, next);
}
\end{minted}

Add the defaults, run the modifiers, then fold the list backwards into a chain
of delegates. All of that happens once, while the schema is being built, and
chapter~\ref{ch:hotchocolate-quickly} established that this now happens at
startup rather than on the first query. A request does not walk a list of
thirteen objects; it enters one delegate that already has the other twelve
closed over inside it.

\index{request pipeline!ConcurrencyGate}%
That is not a detail for its own sake, and the concurrency gate shows why.
Mosaic configures no gate, so at composition time the middleware factory looks
for one, does not find it, and returns the next delegate unwrapped:

\begin{minted}{csharp}
if (gate is null or { IsEnabled: false })
{
    // No gate is configured - skip the middleware entirely so there is
    // no per-request overhead on the hot path.
    return context => next(context);
}
\end{minted}

The comment is theirs, and I repeat it because it is the clearest statement of
the design in the whole file: a middleware with no work to do removes itself
when the pipeline is composed rather than checking on every request. That option
only exists because composition is separate from execution.

Figure~\ref{fig:ch03-pipeline} lays the thirteen out in order, with the places a
request can leave early marked on it. Three of the exits matter more than the
rest, and the sections that follow take them one at a time.

\begin{figure}[htbp]
  \centering
  % The thirteen middleware, and the three places a request can leave early.
% Referenced from chapter 03. Bare tikzpicture: the figure environment, caption
% and label live at the call site. Uses only arrows.meta and positioning.
\begin{tikzpicture}[
    font=\small,
    mw/.style={draw, rounded corners=2pt, minimum width=52mm,
               minimum height=7mm, align=center, font=\scriptsize},
    gate/.style={draw, thick, rounded corners=2pt, minimum width=52mm,
                 minimum height=7mm, align=center, font=\scriptsize},
    outside/.style={draw, dashed, rounded corners=2pt, minimum width=52mm,
                    minimum height=9mm, align=center, font=\scriptsize},
    flow/.style={-{Stealth[length=2mm]}, thick},
    exit/.style={-{Stealth[length=2mm]}, thick, dashed},
    note/.style={font=\scriptsize\itshape, align=left}
  ]

  % The transport, which is not part of the pipeline.
  \node[outside] (http) at (0,0.9) {the ASP.NET Core transport:\\reads the body and \emph{parses the document}};

  \node[mw]   (m1)  at (0,-0.4) {1. Instrumentation};
  \node[mw]   (m2)  at (0,-1.2) {2. Exception};
  \node[mw]   (m3)  at (0,-2.0) {3. Timeout};
  \node[gate] (m4)  at (0,-2.8) {4. DocumentCache};
  \node[mw]   (m5)  at (0,-3.6) {5. DocumentParser};
  \node[gate] (m6)  at (0,-4.4) {6. DocumentValidation};
  \node[gate] (m7)  at (0,-5.2) {7. CostAnalyzer};
  \node[gate] (m8)  at (0,-6.0) {8. OperationCache};
  \node[mw]   (m9)  at (0,-6.8) {9. OperationResolver (compiles)};
  \node[mw]   (m10) at (0,-7.6) {10. SkipWarmupExecution};
  \node[mw]   (m11) at (0,-8.4) {11. OperationVariableCoercion};
  \node[gate] (m12) at (0,-9.2) {12. ConcurrencyGate};
  \node[mw]   (m13) at (0,-10.0) {13. OperationExecution};

  \node[outside] (res) at (0,-11.3) {your resolvers};

  \draw[flow] (http) -- (m1);
  \foreach \a/\b in {m1/m2, m2/m3, m3/m4, m4/m5, m5/m6, m6/m7, m7/m8,
                     m8/m9, m9/m10, m10/m11, m11/m12, m12/m13}{%
    \draw[flow] (\a) -- (\b);
  }
  \draw[flow] (m13) -- (res);

  % The three ways out.
  \draw[exit] (http.east) -- (3.3,0.9);
  \node[note, anchor=west] at (3.4,0.9) {will not parse:\\400, and no request};

  \draw[exit] (m6.east) -- (3.3,-4.4);
  \node[note, anchor=west] at (3.4,-4.4) {fails validation:\\400, nothing executed};

  \draw[exit] (m7.east) -- (3.3,-5.2);
  \node[note, anchor=west] at (3.4,-5.2) {\texttt{GraphQL-Cost: validate}:\\cost only, no data};

  % The thick boxes are the four that can change what the rest of the pipeline
  % has left to do.
  \node[note, anchor=east] at (-3.2,-2.8) {a hit here skips\\5 and 6};
  \node[note, anchor=east] at (-3.2,-4.4) {rejects here, or\\lets the rest run};
  \node[note, anchor=east] at (-3.2,-6.0) {a hit here skips 9};
  \node[note, anchor=east] at (-3.2,-9.2) {removes itself unless\\a gate is configured};

\end{tikzpicture}

  \caption{The thirteen, and the three ways out. A request that does not parse
  never reaches the pipeline at all, because the transport parses it first. A
  request that fails validation stops at six. A request carrying
  \code{GraphQL-Cost: validate} stops at seven with a cost report and no data.
  Only the ones that survive all three reach a resolver.}
  \label{fig:ch03-pipeline}
\end{figure}

\subsection*{An ordering mistake worth repeating}

\index{request pipeline!modifier order}%
The first version of the code that printed that list reported twelve middleware,
not thirteen, and I believed it long enough to write it down.

The reporter works by adding a pipeline modifier, which is the same mechanism
\code{after:} uses. Modifiers run in registration order, and I had registered
mine before the call to \code{AddGraphQL()}. So mine ran first, printed the
pipeline as it stood, and only afterwards did the cost analyzer's modifier
insert the thing I had just finished counting. Moving one line down the file
fixed it.

I would not raise it except that the same trap is waiting for anything else you
insert by name. If your middleware has to sit after a middleware that some
package adds through a modifier, the order in which the two of you were
registered decides whether the anchor is there to find. The failure is loud when
the anchor is missing, and silent when you are merely observing.

So: thirteen middleware, composed once, one of which removes itself entirely
when it has nothing to do. The next four sections walk the ones that stayed,
starting with the one that does its work somewhere else entirely.
\index{request pipeline|)}
